
\section{Decision-analytic evaluation}
\label{sec:dca}

An increment in AUC is not a quantity a service desk can act on. This section
reads the same surface in net benefit, which is: the number of correctly
nominated cases per thousand arrivals, net of the false nominations, at a
declared exchange rate between the two.

\subsection{The operating point is a cost ratio, not a knob}

At threshold $\theta$ the net benefit of a rule that nominates when
$\hat p \geq \theta$ is
\[
\mathrm{NB}(\theta) \;=\; \frac{\mathrm{TP}(\theta)}{n}
\;-\; \frac{\mathrm{FP}(\theta)}{n}\cdot\frac{\theta}{1-\theta},
\]
so $(1-\theta)/\theta$ is the number of false nominations the desk is willing
to accept for one true one. Choosing $\theta$ is choosing an operational cost
ratio, and it belongs to the analyst's declaration rather than to a grid
search. We report the whole grid --- \nGrid\ points from \gridLo\ to
\gridHi\ --- and name three points a desk might occupy: $\theta =
\thetaCheap$, or \ratioCheap\ false nominations accepted per true one, where a
review is cheap and the desk nominates liberally; $\thetaBase$, or
\ratioBase, where a review and a miss cost about the same; and $\thetaDear$,
or \ratioDear, where review capacity is scarce. A reader whose desk sits at a
ratio not in that list reads their own point off Figure~\ref{fig:dca}: the
operating point is the reader's to choose, so the analyst reports the curve
rather than a point on it.

\subsection{Calibration comes first, or the threshold means nothing}
\label{sec:calibfirst}

A threshold is a statement about probability, so a model whose probabilities
are wrong is being read at a threshold it does not have. Every decision curve
here is therefore computed on \emph{calibrated} scores, with the calibrator
fitted inside the training half and rebuilt inside every bootstrap refit, and
a registered rule excludes an arm whose calibration slope leaves
\calSlopeWindow.

The window is a declared convention and not a derived bound, and the exclusion
is deliberately not repaired by recalibrating the failing arm: recalibration is
already a level of the pipeline axis, so repairing an arm with it would replace
one specification by another and report the result under the first one's name.
Supplement~\ref{app:dcadetail} gives the argument, the rule, what it removes
and what the curves look like without it.

\subsection{Simultaneous bands change what can be said}
\label{sec:dcabands}

The decision-curve family is declared separately from the scalar surface and
carries its own max-$t$ critical value, because net benefit and a scalar
instrument are not the same scale. Its critical value is the empirical
quantile of Section~\ref{sec:simbands}, taken over the \emph{admissible}
family: a cell whose net-benefit difference is exactly zero in at least nine
draws of ten is zero by arithmetic, not by estimation, and is excluded before
the quantile is taken because it is not a statement about the register while
its vanishing spread inflates the value every other cell is judged by
(Section~\ref{sec:simband}). On this estate's family that rule excludes
\nDcaCellsExcludedDegenerate\ of \nDcaFamily\ cells and gives a critical
value of \dcaQEmp\ [\dcaQEmpLo, \dcaQEmpHi], against \dcaQEmpAll\ with those
cells kept and \dcaQ\ from the multiplier approximation.
Section~\ref{sec:simband} measures this estimator at \covEmpHeavyCaseAtDesign\
on a family matched to this one and the multiplier at
\covMultHeavyCaseAtDesign, so the band below is the better of the two and is
still about a point short of nominal on the matched family, which is stated
here rather than in a footnote.

Across the \nDcGrid\ operating points of this estate's curve \textbf{the
pointwise interval declares the increment resolvably positive at
\nBeneficialPointwise\ of them and the simultaneous band at
\nBeneficialSimultaneousEmp}, the lowest of the second at $\theta =
\thetaBeneficialEmpMin$; the harmful counts are \nHarmfulPointwise\ against
\nHarmfulSimultaneousEmp. All are built from the same draws, so what differs
is the family and the level it is bought at: those points are the price of a
statement about a \emph{range} of thresholds rather than about each one
separately, and a paper that reads a decision curve across its range while
pricing it pointwise has claimed the first and paid for the second. The
multiplier approximation would have declared \nBeneficialSimultaneous\ of the
\nDcGrid\ resolved, at a median band width of \dcBandWidth\ against the
reported band's \dcBandWidthEmp\ and the pointwise interval's
\dcPointWidth; the order-statistic upper end of the reported critical value
leaves \nBeneficialSimultaneousEmpHi, and keeping the excluded cells leaves
\nBeneficialSimultaneousEmpAll, so the exclusion rule does not carry the
result. Supplement~\ref{app:dcabandsfull} prints the whole ladder,
including the measured widening of the multiplier value that an earlier
construction used, and Figure~\ref{fig:dca} and Table~\ref{tab:dca} the
curve. The pointwise interval is not widened by a family-wise factor; what
it covers is a separate measurement, and the $(n,K)$ calibration of
Section~\ref{sec:regions} is not applied to this family either.

\subsection{The decision, in units a service desk uses}
\label{sec:nbregret}

The comparison is run in net benefit per thousand cases, over a declared
distribution of exchange rates a desk plausibly holds
(\deskQlo\ to \deskQhi, centred \deskQmid), and only on models that pass the
registered calibration rule --- which excludes \nModelsDcaExcluded\ of
\nModelsDca\ arms and \nPairsDcaExcluded\ of \nPairs\ pairs.
\textbf{That count is itself the strongest operational finding here}: a
decision curve read off a model whose probabilities are wrong is read at a
threshold the model does not have.

\textbf{And the rules separate on most pairs, with the excess falling on the
conservative rule rather than the one-number rule.} The one-number rule's
median excess is \excessOneNumberDesk\ per thousand against the conservative
rule's \excessUniformDesk, and the cross-pair means are
\meanExcessOneNumberDesk\ and \meanExcessUniformDesk: what the medians hide
is a tail, and it is the conservative rule's. The four rules do not all lose
the same amount on \nSeparatingDesk\ of the \nPairsDesk\ pairs the
calibration rule leaves, and all four are identical on the other
\nAllRulesTieDesk. Where they do separate, the one-number rule is strictly
the worst on \nOneNumberWorstDesk\ and tied for worst on
\nOneNumberTiedWorstDesk, at a maximum of \maxExcessOneNumberDesk\ per
thousand.
\textbf{The case against one-number reporting does not rest on this design.}
Larger than anything the choice of collapse rule costs a one-number report is
the decision time: the register promises \promisedPerThousand\ per thousand
at the earlier moment and delivers \deliveredPerThousand\ at the later one, a
shortfall of \shortfallPerThousand. Supplement~\ref{app:deskfull} carries the
rules, the exchange-rate distribution and the exclusions;
Tables~\ref{tab:desk} and~\ref{tab:rules} the numbers.
